% Technical Background, section 2. CAS-Net subsection and methodology/02_segmentation.tex dropped
% (user, 2026-10-21): lumen and centerline now come from the in-house tracer (sec:method-tracer).
% TODO: keep this U-Net section only if the tracer's background needs it; otherwise cut it.
% Model: Aida 3.2 (U-Net, section body) and 3.2.1 (nnU-Net, the subsection one step deeper).
% Architecture sequence (guide section 3.3): origin, structure, key characteristic, why it works,
% standing, tool. Convolution and gradient-based training are assumed, as in Aida; the soft Dice
% loss goes to the Methodology next to Dice.
% Cite: ronneberger2015unet, isensee2021nnunet, wasserthal2023totalsegmentator.
% thesis/pipeline/ (the CAS-Net chapter) was deleted; its prose is in git history (week6).
% FIGURE WANTED: own 3D U-Net schematic (academic-figures skill), or Ronneberger's with
%   credit in figures/external/CREDITS.md. The chapter's only figure, as in Aida.

\section{3D Segmentation: U-Net}
\label{sec:tech-unet}
% P1 definition opener and origin: U-Net is a convolutional network for biomedical segmentation,
%    introduced by Ronneberger et al. in 2015, trained with few labelled images.
% P2 structure: encoder (contracting path) reduces resolution and widens channels, in 3D for volumes.
% P3 decoder (expanding path) restores resolution to a voxel-wise class map.
% P4 key characteristic: skip connections carry encoder features to the matching decoder level.
% P5 why it works, tied to the object: thin branching vessels need both context and voxel precision.
